\section{Fase 4. Desarrollo del Director del juego}
\label{sec:F4Intro}

Hasta ahora disponemos de un director del juego básico, que lanza los
dos motores necesarios y poco más. Esta fase saca del código todos los
parámetros que gobiernan el ritmo de la partida y los convierte en
propiedades del cartucho, editables desde el diseñador gráfico. La
motivación es doble: por un lado, el simulador de PC y el futuro
sistema sobre ESP32 comparten toda la lógica del motor de estados,
que es determinista en \textit{ticks}, pero sus procesadores trabajan
a velocidades muy distintas y pueden exigir pausas de reloj diferentes
para ofrecer el mismo \textit{feel} de juego; por otro, la fase
siguiente (el generador del cartucho compacto) necesita que
\textbf{todo} lo configurable viaje en el GWI, incluido lo que hasta
ahora solo vivía en el setup o en constantes del código.

Cerraremos en esta fase dos temas:

\begin{itemize}
	\item Ajustar el Director del juego: todos los parámetros que
	gobiernan el ritmo y la dificultad de la partida dejan de ser
	constantes del código o del \emph{setup} y pasan a ser propiedades
	del cartucho, editables desde el diseñador gráfico y visibles para el
	motor al instante, sin guardar ni recargar. En concreto: una
	velocidad de reloj por plataforma, una aceleración que reduce la
	pausa con el tiempo y una política de lanzamiento de rutas con
	densidad creciente. Con esto queda cerrado en el PC todo lo
	configurable del juego, requisito de la fase siguiente.
	\item Como objetivo secundario, cerrar todos los temas pendientes
	para la generación del cartucho intermedio, que debe llevar toda la
	información necesaria para generar el cartucho compacto.
\end{itemize}

Al igual que en las fases anteriores, los parámetros pueden llegar
opcionalmente del archivo de setup, pero la fuente de verdad durante
la autoría es siempre el diseñador gráfico: cualquier edición en la
ficha del cartucho se sincroniza con el motor al instante, sin
guardar ni recargar.

\section{Primera iteración: Análisis de los cambios necesarios}
\label{sec:F4I1}

\subsection{Diferentes velocidades por plataforma}

El reloj del Director encadena sus \textit{ticks} con
\texttt{setTimeout} y una pausa fija de 500 milisegundos escrita en el
código (\texttt{PAUSA = 500}). Este valor es bastante correcto para el
simulador en PC\footnote{En mi ordenador, usar 480 es más aproximado
	al G\&W real; con el ritmo en el cartucho, ese ajuste fino pasa a ser
	una propiedad del GWI de cada juego y ya no del código.}, pero no
tiene por qué serlo para el ESP32: la duración real de cada
\textit{tick} (barrido, motor de estados y renderizado) depende del
procesador, y replicar la velocidad subjetiva de la máquina original
puede exigir pausas distintas en cada plataforma. Además, un valor
fijo en el código obliga a recompilar para probar ritmos distintos, y
la fase siguiente necesita que el ritmo sea una propiedad del
cartucho para que el sistema sobre ESP32 lo reciba ya resuelto.

\subsection{Aceleración del juego con el tiempo}

En las máquinas reales la partida no mantiene un ritmo constante: a
medida que avanza el tiempo, el juego se vuelve más difícil. El
simulador mantenía la misma pausa durante toda la partida, de modo
que la dificultad solo crecía por acumulación de desenlaces, no por
velocidad. Falta el mecanismo que reduzca la pausa del reloj con el
tiempo, de forma que la partida acelere sola, sin necesidad de modos
de juego separados.

\subsection{Política de lanzamiento de rutas}

El nacimiento de rutas era una política fija del motor de estados:
solo nacía una ruta nueva en el \textit{tick} en que había finalizado
la vida de la anterior (las animaciones pueden ser simultáneas, pero
no son parte del juego ni son rutas que se lanzan). En las máquinas
reales la densidad de tráfico es parte del diseño de cada juego y
crece con la partida: comienzan cruzando pocas rutas y, a medida que
avanza, van aumentando las rutas simultáneas. Falta, por tanto, que
esta política se pueda redefinir y suba al cartucho, editable desde
el diseñador gráfico.

\subsection{Revisión de campos en el cartucho}

Como hay que añadir todos los temas de configuración general en el
cartucho para su uso por el Director del juego y por el diseñador de
cartuchos compactos, se revisa que incluya toda la información
necesaria para generar el cartucho final.

\section{Segunda iteración: Cambios necesarios}
\label{sec:F4I2}

Esta iteración implementa el análisis anterior. Los tres bloques
(velocidades, aceleración y política de rutas) comparten el mismo
camino de datos, así que se desarrollaron juntos: viajan en el GWI,
cruzan el puente \texttt{ComunListas} y se editan en la ficha del
cartucho del diseñador, con el mismo patrón que ya usaban el fin de
partida y la Secuencia.

\subsection{El GWI ampliado}

El cartucho gana tres campos al nivel superior. Un GWI que declara
todo queda así (fragmento):

\begin{lstlisting}[language=json, caption={Campos del Director en el
		GWI (fragmento).}]
	"pausas":      { "simulador": 480, "externo": 500 },
	"aceleracion": { "cada": 12, "resta": 5, "minimo": 400 },
	"rutas":       { "aumentarCada": 3 }
\end{lstlisting}

\begin{itemize}
	\item \texttt{pausas}: las dos velocidades de reloj, una por
	plataforma. El simulador lee siempre \texttt{simulador}; el firmware
	del ESP32 leerá \texttt{externo}. Si el cartucho no declara el campo,
	el Director mantiene sus 500 ms históricos en ambas plataformas.
	\item \texttt{aceleracion}: cada \texttt{cada} \textit{ticks} de
	partida la pausa vigente se reduce en \texttt{resta} milisegundos,
	sin bajar nunca de \texttt{minimo}. Un solo conjunto de parámetros
	sirve a ambas plataformas: la aceleración se aplica sobre la pausa
	inicial de cada una y el efecto es proporcional. Ausente, la pausa
	permanece fija durante toda la partida.
	\item \texttt{rutas}: \texttt{aumentarCada} es el «Nº para aumentar
	rutas». Ausente, el máximo de simultáneas queda en 1 (la regla
	histórica).
\end{itemize}

\subsection{Los tres contadores de rutas}

La semántica de la política de lanzamiento se apoya en tres
contadores bien separados, porque confundirlos fue el principal
riesgo de diseño de la fase:

\begin{itemize}
	\item \textbf{Total de rutas lanzadas} (en el Director): nace en 0 y
	sube con cada lanzamiento; \textbf{nunca} baja con las muertes. Es el
	contador de progresión de la partida: la densidad crece aunque las
	rutas mueran.
	\item \textbf{Máximo de rutas simultáneas} (en el Director): nace en
	1 y sube en 1 cada vez que el total alcanza un múltiplo de
	\texttt{aumentarCada}. Su tope real es el número de rutas definidas
	en el cartucho (solo hay tantas rutas como inicios), aplicado en el
	motor al elegir candidato.
	\item \textbf{Rutas vivas} (en el motor de estados, que es el dueño
	de la pantalla): cuántas hay en pantalla ahora mismo. En cada
	\textit{tick}, si las vivas son menos que el máximo de simultáneas,
	se lanza una nueva (una por tick, para que el nacimiento sea
	visible). La regla histórica (nacer cuando no queda ninguna viva) es
	el caso particular de máximo igual a 1.
\end{itemize}

Al repartirlos así, cada contador vive con su dueño natural: el
Director lleva la progresión (lanza y autoriza) y el motor lleva la
pantalla (cuenta lo que hay encendido). La comunicación usa dos
funciones nuevas de la API: el motor llama a
\texttt{JuegoDirector.rutaLanzada()} al encender el inicio de una
ruta, y consulta \texttt{JuegoDirector.maxRutasSimultaneas()} antes
de decidir; el motor expone además
\texttt{estadoRutas()}, que devuelve \texttt{\{vivas,
	maxSimultaneas\}} para las pruebas y el diagnóstico.

\subsection{El camino de los datos: el puente}

\texttt{ComunListas} gana tres pares de accesos, siguiendo el patrón
del fin de partida y de la Secuencia:
